% Chapter 5

\chapter{SOCIAL IMPACT AND SUSTAINABILITY}\label{cha:social}

\section{SOCIAL IMPACT}

Drone delivery helps most the people who live far from a hospital or are cut off by an
emergency. Rural healthcare logistics considerations would be in cases where there was a
need for blood or an antivenom in a clinic, a town surrounded by floods, or regular
delivery of samples from a hospital to a laboratory in one day. Similar services already exist, such as Zipline drones in Rwanda~\cite{r16} and hospital samples flights in Switzerland~\cite{r12}. The suggested method adds a way for
each flight to adapt to the particular cargo, where urgent antivenom would fly through
the shortest path while fragile cargo would use another path. This advantage is important since there might be a situation where no specialists in
planning will be available. The launch of a drone from a district hospital or a relief
camp will be conducted either by a nurse or volunteers instead of a specialist in
planning. This person needs only to provide four simple answers about the cargo
characteristics (frangible, weight, urgency, and importance), after which all other
decisions will be made automatically.

Underneath, there are people who do not want the drone to fly above their heads. Every
edge cost considers aspects such as population, restricted areas, water distance, and
building distance. All nodes that exceed the statutory limit are disregarded, and
adverse weather conditions make it impossible for flights if there are storms, rainfall,
poor visibility, and high winds. There is no way that the planner will go through a
crowded market just to save a few seconds.

All decisions made at Tier~2 level are saved along with the corresponding tags, weights,
route, and any HOLD reason. The fragility tag
is a precaution and has not been shown to protect anything~\cite{r13,r14,r15}. Until
proper risk weights are set, the derived risk values will act as relative rankings of
the routes rather than probabilities of harm. Phase~II outdoor flights will follow the
Drone Rules, 2021~\cite{r21}.

\section{SUSTAINABILITY}

For the short-distance flights of light payloads, an electric-powered drone may be
substituted for van or motorcycle transport without exhaust emission during travel. The
planner helps in two major ways. Firstly, its energy goal considers the wind direction,
which alone changed the energy of the cheapest route by 27.5\%, implying that headwinds
may be avoided. Secondly, the payload weight label gives more weight to energy criteria
for heavy payloads. Flights whose battery energy is exhausted to less than the 20\%
reserve level are excluded, thereby assuring both the success of the flight and battery
longevity.

The computational demand is low as well. Tier~1 runs on an ordinary laptop CPU with no
GPU or cloud service, once per region and weather window. Tier~2 takes a fraction of a
millisecond and can run on a Raspberry Pi on board. Everything the system depends on is open data or open-source software, so there are no licence fees.
